% Source: https://artofproblemsolving.com/community/c3h3181916_2021_abmc_online_november_contest__actonboxborough_math_competition
% Stable identity: aops-topic-3181916-q13
Scale the inradius to 1 and put $B=(-\sqrt3,0),C=(\sqrt3,0),A=(0,3)$; the incenter is $I=(0,1)$. Angle $BPC$ is acute exactly when $P$ lies outside the disk with diameter $BC$, whose radius is $\sqrt3$ and centre is $(0,0)$. The circles intersect at $(\pm\dfrac{\sqrt3}{2},\dfrac{3}{2})$. Above their common chord, the incircle segment has area $\dfrac{\pi}{3}-\dfrac{\sqrt3}{4}$, while the diameter-circle segment has area $\dfrac{\pi}{2}-\dfrac{3\sqrt3}{4}$. Their difference is $\dfrac{\sqrt3}{2}-\dfrac{\pi}{6}$. Divide by incircle area $\pi$ to get $\dfrac{3\sqrt3-\pi}{6\pi}$, giving $\boxed{13}$.
\begin{asy}
usepackage("amsmath");
size(7cm);
pair A=(0,3),B=(-sqrt(3),0),C=(sqrt(3),0),Ii=(0,1);draw(A--B--C--cycle);draw(circle(Ii,1),blue);draw(circle((0,0),sqrt(3)),red);draw((-sqrt(3)/2,1.5)--(sqrt(3)/2,1.5),dashed);label("$A$",A,N);label("$B$",B,SW);label("$C$",C,SE);label("$I$",Ii,W);
\end{asy}
